\chapter{その他}
\section{集合論}
\subsection{濃度}

\begin{egprob}
$A$を次の性質を満たす関数$f:\mb{R}\to \mb{R}$全体のなす集合とする。
\ben[label=(\roman*)]
\item $f(1) = 1$
\item $f(x+y) = f(x) + f(y)$
\een
\ben
\item $f\in A$のとき, $f(x) = x$ ($x\in \mb{Q}$)であることを示せ。
\item $r\in \mb{R}\setminus{\mb{Q}}$を固定する。$f(r) = 0$であるような$f\in A$が存在することを示せ。
\item $A$の部分集合$\{ f\in A\mid \text{$f$は$\mb{R}$のある点で不連続} \}$の濃度を求めよ。
\een
\end{egprob}
\begin{egans}
(1): $f(0) = 0$, $f(n) = n = -f(-n)$ ($n\in \mb{N}$) は容易。$x=m/n$のとき$m=f(m) = f(nx) = nf(x)$だから$x\in \mb{Q}$のときもよい。\qed

(2): $\{ 1,r\}$は$\mb{Q}$上一次独立である。$\mb{R}$は自然に$\mb{Q}$ベクトル空間の構造を持つ。よって$\{ 1,r\}$はある基底$\{ e_i\}_{i\in I}$に拡張される。これを用いて$x=\sum_{i\in I} a_{i}(x)e_i$となる$a_i(x)\in \mb{Q}$が一意に取れる。ただしほとんどすべての$i\in I$に対して$a_i(x) = 0$である。$e_t = 1$, $e_s = r$であるとする。
このとき, $f:\mb{R}\to\mb{R}$を
\[f(x) = a_{t}(x)e_t = a_{t}(x)\]
と定めるとよい。$\{ e_i \}_{i\in I}$が基底だから, $a_{t}(x+y) = a_{t}(x) + a_{t}(y)$となる。$a_{t}(1) = 1$だから$f(1) = 1$であり, $a_{t}(r) = 0$なので$f(r) = 0$である。この$f$でよい。

(3): $f\in A, i\in I$に対して, $f(e_i) = f_i$とする。条件から$f_t = 1$である。こうして$f\in A$から実数の族$(f_i)_{i\in I\setminus{\{ t \}}}\in \mb{R}^{I-\setminus{\{ t \}}}$が得られる。逆に, $I' = I -\setminus{\{ t \}}$とし, 任意の$(k_i)\in \mb{R}^{I'}$に対し, $f(e_i) = k_i$となる$f\in A$を
\[f(x) = \sum_{i\in I'} a_{i}(x)k_{i}\]
で定めることができる(和は本質的に有限であるからwell-defined)。こうして$A$と$\mb{R}^{I'}$の間に全単射がある。$\mb{Q}$ベクトル空間として$\oplus_{i\in I}\mb{Q} \cong \mb{R}$であり, 左辺の濃度は$|\mb{Q}|^{|I|}=2^{|I|}$だから$\mb{Q}$が可算で$\mb{R}$が連続濃度であることから$|I'| = |\mb{N}|$である。よって$I'$も可算で$|\mb{R}^{I'}| = |\mb{R}|=\aleph$であるから$|A| = |\mb{R}|$. $f\in A$は$f|_{\mb{Q}} = \mr{id}_{\mb{Q}}$を満たしていたから, $f$がすべての点で連続であるとするなら$f=\mr{id}_{\mb{R}}$でなければならない。よって$A-\{ \mr{id}_{\mb{R}}\}$の濃度を求めればよいが, これも$|\mb{R}|$に等しい。よって濃度は$\aleph$.

\end{egans}
\subsubsection{HINTS}
\begin{hint}
\chatgpthint{$\mb{R}$を$\mb{Q}$ベクトル空間とみなし、Hamel基底を用います。}
\end{hint}




